\chapter{Divisors and Line Bundles}

\section{Weil Divisors}

\begin{definition}
Let \(X\) be an integral noetherian scheme.  A \emph{prime divisor} on \(X\) is
an integral closed subscheme of codimension one, or equivalently a point
\(\eta\in X\) whose closure has codimension one.  A \emph{Weil divisor} is a
finite formal sum
\[
    D=\sum_\eta n_\eta[\eta],
\]
where \(n_\eta\in\ZZ\) and \(\eta\) runs over codimension-one points.  The group
of Weil divisors is denoted \(\Div(X)\).  A divisor is \emph{effective}, written
\(D\ge0\), if all \(n_\eta\ge0\).
\end{definition}

\begin{definition}
Assume that \(X\) is normal.  For \(f\in k(X)^\times\), the
\emph{principal divisor} of \(f\) is
\[
    \operatorname{div}(f)=
    \sum_{\codim(\eta)=1}\ord_\eta(f)[\eta],
\]
where \(\ord_\eta\) is the discrete valuation of the DVR \(\OO_{X,\eta}\).
\end{definition}

\begin{proposition}
For \(f,g\in k(X)^\times\),
\[
    \operatorname{div}(fg)=\operatorname{div}(f)+\operatorname{div}(g).
\]
Thus principal divisors form a subgroup of \(\Div(X)\).
\end{proposition}

\begin{proof}
For every codimension-one point \(\eta\), the valuation
\(\ord_\eta:k(X)^\times\to\ZZ\) is a group homomorphism.  Therefore
\[
    \ord_\eta(fg)=\ord_\eta(f)+\ord_\eta(g)
\]
for each \(\eta\).  Since equality of Weil divisors is equality of all
coefficients, the formula follows.
\end{proof}

\begin{definition}
The \emph{class group} of a normal noetherian integral scheme \(X\) is
\[
    \Cl(X)=\Div(X)/\{\text{principal divisors}\}.
\]
\end{definition}

\section{Cartier Divisors}

\begin{definition}
Let \(X\) be a scheme.  The sheaf of total quotient rings \(\mathcal K_X\) is
defined on an affine open \(U=\Spec A\) by
\[
    \mathcal K_X(U)=S^{-1}A,
\]
where \(S\) is the multiplicative set of non-zero-divisors in \(A\).  A
\emph{Cartier divisor} is a global section of
\(\mathcal K_X^\times/\OO_X^\times\).  Equivalently, it is given by an open
cover \(X=\bigcup_i U_i\) and elements \(f_i\in\mathcal K_X^\times(U_i)\) such
that \(f_i/f_j\in\OO_X^\times(U_i\cap U_j)\).
\end{definition}

\begin{definition}
A Cartier divisor represented by \((U_i,f_i)\) is \emph{principal} if it is
represented by one global element \(f\in\mathcal K_X^\times(X)\).  It is
\emph{effective} if it can be represented locally by non-zero-divisors
\(f_i\in\OO_X(U_i)\).
\end{definition}

\begin{proposition}
Cartier divisors form an abelian group \(\CaDiv(X)\), with addition represented
locally by multiplication of the local equations.
\end{proposition}

\begin{proof}
If \(D\) is represented by \((U_i,f_i)\) and \(E\) by \((V_j,g_j)\), then on the
cover \(U_i\cap V_j\) define \(D+E\) by \(f_i g_j\).  On overlaps,
\[
    \frac{f_i g_j}{f_{i'}g_{j'}}
    =
    \frac{f_i}{f_{i'}}\frac{g_j}{g_{j'}}
\]
is a unit, so the data define a Cartier divisor.  The zero element is
represented by \(1\), and the inverse of \((U_i,f_i)\) is \((U_i,f_i^{-1})\).
Associativity and commutativity follow from multiplication in
\(\mathcal K_X^\times\).
\end{proof}

\section{The Sheaf \texorpdfstring{\(\OO_X(D)\)}{O(D)}}

\begin{definition}
Let \(D\) be a Cartier divisor represented by local equations
\((U_i,f_i)\).  Define an invertible subsheaf \(\OO_X(D)\subset \mathcal K_X\)
by
\[
    \OO_X(D)|_{U_i}=f_i^{-1}\OO_{U_i}.
\]
\end{definition}

\begin{proposition}
The sheaf \(\OO_X(D)\) is well defined and invertible.  Moreover
\[
    \OO_X(D+E)\simeq \OO_X(D)\tensor_{\OO_X}\OO_X(E).
\]
\end{proposition}

\begin{proof}
On \(U_i\cap U_j\), the quotient \(f_i/f_j\) is a unit.  Hence
\[
    f_i^{-1}\OO_{U_i\cap U_j}
    =
    f_j^{-1}\OO_{U_i\cap U_j}
\]
as subsheaves of \(\mathcal K_X\).  The local definitions therefore glue to a
well-defined sheaf.  On \(U_i\), multiplication by \(f_i\) identifies
\(\OO_X(D)|_{U_i}\) with \(\OO_{U_i}\), so \(\OO_X(D)\) is invertible.

If \(E\) is represented by \((V_j,g_j)\), then \(D+E\) is represented on
\(U_i\cap V_j\) by \(f_i g_j\).  On this open set,
\[
    \OO_X(D+E)=(f_i g_j)^{-1}\OO
    \simeq f_i^{-1}\OO\tensor_\OO g_j^{-1}\OO.
\]
These local multiplication maps are compatible on overlaps and glue to the
displayed isomorphism.
\end{proof}

\begin{proposition}
If \(D=\operatorname{div}(f)\) is principal, then
\[
    \OO_X(D)\simeq \OO_X.
\]
\end{proposition}

\begin{proof}
The divisor is represented globally by \(f\in\mathcal K_X^\times(X)\).  Hence
\(\OO_X(D)=f^{-1}\OO_X\).  Multiplication by \(f\) gives an isomorphism
\[
    f^{-1}\OO_X\longrightarrow \OO_X
\]
with inverse multiplication by \(f^{-1}\).
\end{proof}

\begin{proposition}
Let \(X\) be an integral scheme and \(D\) a Cartier divisor.  Then
\[
    H^0(X,\OO_X(D))
    =
    \{h\in k(X): \operatorname{div}(h)+D\ge0\},
\]
where the right side is interpreted locally by the Cartier equations of \(D\).
\end{proposition}

\begin{proof}
Let \(D\) be represented on \(U_i\) by \(f_i\).  A rational function \(h\) is a
section of \(\OO_X(D)\) precisely when \(h|_{U_i}\in f_i^{-1}\OO_X(U_i)\) for
all \(i\), equivalently \(hf_i\in\OO_X(U_i)\) for all \(i\).  At each
codimension-one point \(\eta\), this condition says
\[
    \ord_\eta(h)+\ord_\eta(D)\ge0.
\]
That is exactly \(\operatorname{div}(h)+D\ge0\).
\end{proof}

\section{Cartier and Weil Divisors on Regular Schemes}

\begin{proposition}
Let \(X\) be a regular integral noetherian scheme.  Every Cartier divisor
defines a Weil divisor by
\[
    D=(U_i,f_i)\longmapsto
    \sum_{\codim(\eta)=1}\ord_\eta(f_i)[\eta],
\]
where \(i\) is chosen with \(\eta\in U_i\).  This is independent of \(i\).
\end{proposition}

\begin{proof}
If \(\eta\in U_i\cap U_j\), then \(f_i/f_j\) is a unit in
\(\OO_{X,\eta}\).  Since \(\OO_{X,\eta}\) is a DVR for a regular
codimension-one point, units have valuation zero.  Therefore
\(\ord_\eta(f_i)=\ord_\eta(f_j)\).  Only finitely many coefficients are
non-zero because on each affine open a non-zero rational function has zeros
and poles along only finitely many height-one primes in a noetherian domain.
\end{proof}

\begin{proposition}
If \(X\) is a regular integral noetherian scheme, then every Weil divisor is
Cartier locally in codimension one.  If \(X\) is a regular integral curve, the
map from Cartier divisors to Weil divisors is an isomorphism.
\end{proposition}

\begin{proof}
Let \(X\) be a regular curve.  For a closed point \(P\), the local ring
\(\OO_{X,P}\) is a DVR, so its maximal ideal is generated by a uniformizer
\(t_P\).  A Weil divisor \(D=\sum_P n_P[P]\) has finite support.  Around each
point \(P\) in the support, the local equation \(t_P^{n_P}\) defines the
required Cartier divisor.  Away from the support, use the local equation \(1\).
On overlaps, the ratios are units because they have valuation zero at every
point of a regular one-dimensional local ring.  Thus \(D\) is Cartier.

The previous proposition gives the inverse map from Cartier divisors to Weil
divisors, and the two constructions are inverse because both are determined by
the valuation at each closed point.
\end{proof}

\section{Degree on Curves}

\begin{definition}
Let \(X\) be an integral proper curve over a field \(k\).  For a Weil divisor
\[
    D=\sum_P n_P[P]
\]
with \(P\) running over closed points, define
\[
    \deg D=\sum_P n_P[\kappa(P):k].
\]
\end{definition}

\begin{proposition}
The degree map \(\deg:\Div(X)\to\ZZ\) is a group homomorphism.
\end{proposition}

\begin{proof}
If \(D=\sum_P n_P[P]\) and \(E=\sum_P m_P[P]\), then
\[
    \deg(D+E)
    =
    \sum_P (n_P+m_P)[\kappa(P):k]
    =
    \deg D+\deg E.
\]
The sums are finite because divisors have finite support.
\end{proof}

\begin{theorem}
Let \(X\) be a smooth projective integral curve over \(k\).  For every
\(f\in k(X)^\times\),
\[
    \deg\operatorname{div}(f)=0.
\]
\end{theorem}

\begin{proof}
The rational function \(f\) defines a non-constant morphism
\[
    \phi:X\longrightarrow \PP^1_k
\]
when \(f\notin k^\times\); if \(f\in k^\times\), the divisor is zero.  Since
both curves are projective, \(\phi\) is proper, and since the induced extension
of function fields is finite, \(\phi\) is a finite morphism.  The divisor of
the coordinate \(t\) on \(\PP^1\) is \([0]-[\infty]\).  Pulling back Cartier
divisors gives
\[
    \operatorname{div}(f)=\phi^*[0]-\phi^*[\infty].
\]
For a finite morphism of integral proper curves, the degree of the pullback of
a closed point \(Q\) is
\[
    \deg(\phi^*[Q])
    =
    [k(X):k(t)]\,[\kappa(Q):k],
\]
because the tensor product
\(\OO_{\PP^1,Q}\)-length of \(\OO_{X,\phi^{-1}Q}\) equals the degree of the
function field extension.  The points \(0\) and \(\infty\) are \(k\)-rational,
so the two pullbacks have the same degree.  Hence
\(\deg\operatorname{div}(f)=0\).
\end{proof}

\begin{definition}
Let \(X\) be a smooth projective integral curve.  A \emph{canonical divisor} is
any divisor \(K\) such that
\[
    \OO_X(K)\simeq \omega_X.
\]
Such a divisor exists because every invertible sheaf on a regular integral
curve is obtained from a Cartier divisor.
\end{definition}

\section{Picard Group and Cartier Divisors}

\begin{proposition}
For any scheme \(X\), the assignment
\[
    D\longmapsto \OO_X(D)
\]
defines a group homomorphism
\[
    \CaDiv(X)\longrightarrow \Pic(X).
\]
Principal Cartier divisors map to the trivial class.
\end{proposition}

\begin{proof}
The tensor product formula
\[
    \OO_X(D+E)\simeq\OO_X(D)\tensor\OO_X(E)
\]
shows that the assignment is a homomorphism.  If \(D=\operatorname{div}(f)\) is
principal, multiplication by \(f\) gives \(\OO_X(D)\simeq\OO_X\), so principal
Cartier divisors map to the identity in \(\Pic(X)\).
\end{proof}

\begin{proposition}
Let \(X\) be an integral regular noetherian scheme.  Then every invertible
subsheaf \(\LL\subset\mathcal K_X\) is of the form \(\OO_X(D)\) for a Cartier
divisor \(D\).  If \(X\) is a regular integral curve, every invertible sheaf is
isomorphic to \(\OO_X(D)\) for some divisor \(D\).
\end{proposition}

\begin{proof}
Choose an open cover \(U_i\) on which \(\LL\) is generated by a rational
section \(s_i\in\mathcal K_X^\times(U_i)\).  On \(U_i\cap U_j\), the quotient
\(s_i/s_j\) is a unit because both generate the same invertible subsheaf.
Therefore the \(s_i^{-1}\) define a Cartier divisor \(D\), and
\(\LL=\OO_X(D)\).  On a regular integral curve, any invertible sheaf admits a
non-zero rational section: trivialize it at the generic point, then view local
generators as rational multiples of this trivialization.  Hence the first part
applies.
\end{proof}

\begin{definition}
Two divisors \(D,E\) on an integral regular scheme are \emph{linearly
equivalent}, written \(D\sim E\), if \(D-E\) is principal.
\end{definition}

\begin{corollary}
On a regular integral curve, the map \(D\mapsto\OO_X(D)\) induces an
isomorphism
\[
    \Div(X)/\sim \;\simeq\; \Pic(X).
\]
\end{corollary}

\begin{proof}
Surjectivity is the preceding proposition.  If \(\OO_X(D)\simeq\OO_X(E)\), then
\(\OO_X(D-E)\simeq\OO_X\).  A trivialization of \(\OO_X(D-E)\) is given by a
non-zero rational function \(f\) satisfying
\[
    \operatorname{div}(f)+D-E=0.
\]
Thus \(D-E\) is principal.  The converse was already shown: principal divisors
map to the trivial line bundle.
\end{proof}

\section{Pullback of Divisors}

\begin{definition}
Let \(f:X\to Y\) be a morphism and let \(D\) be a Cartier divisor on \(Y\)
represented by \((V_i,g_i)\).  If no \(f(X)\) is contained in the support where
the local equation becomes a zero-divisor, the \emph{pullback} \(f^*D\) is the
Cartier divisor represented on \(f^{-1}V_i\) by \(f^\#g_i\).
\end{definition}

\begin{proposition}
Under the preceding hypothesis,
\[
    \OO_X(f^*D)\simeq f^*\OO_Y(D).
\]
\end{proposition}

\begin{proof}
On \(V_i\), \(\OO_Y(D)\) is \(g_i^{-1}\OO_Y\).  Pulling back to
\(f^{-1}V_i\) gives \((f^\#g_i)^{-1}\OO_X\), which is exactly the local
description of \(\OO_X(f^*D)\).  The local identifications agree on overlaps
because the transition functions \(g_i/g_j\) pull back to the transition
functions \(f^\#g_i/f^\#g_j\).
\end{proof}

\begin{proposition}
Let \(f:X\to Y\) be a finite morphism of smooth projective curves.  For a
closed point \(Q\in Y\),
\[
    f^*[Q]=\sum_{P\mapsto Q} e_P [P],
\]
where \(e_P\) is the ramification index at \(P\).  Consequently
\[
    \deg f^*[Q]=(\deg f)[\kappa(Q):k].
\]
\end{proposition}

\begin{proof}
Choose a uniformizer \(u\) at \(Q\).  At \(P\), the pullback \(f^\#u\) has order
\(e_P\) by definition of ramification index, so the coefficient of \(P\) in the
pullback divisor is \(e_P\).  For the degree formula, localize the finite
\(\OO_{Y,Q}\)-module \(f_*\OO_X\) at \(Q\).  Its length after quotienting by
\((u)\) is
\[
    \sum_{P\mapsto Q} e_P[\kappa(P):\kappa(Q)].
\]
This length is also the degree of the finite extension of function fields
\(\deg f\).  Multiplying by \([\kappa(Q):k]\) gives the displayed equality.
\end{proof}

\section{Linear Systems}

\begin{definition}
Let \(X\) be a smooth projective curve and let \(D\) be a divisor.  The
\emph{complete linear system} \(|D|\) is the set of effective divisors linearly
equivalent to \(D\).  Equivalently,
\[
    |D|=\PP(H^0(X,\OO_X(D)))
\]
after identifying a non-zero section with its divisor of zeros plus \(D\).
\end{definition}

\begin{proposition}
Let \(s\in H^0(X,\OO_X(D))\) be non-zero.  Then there is a unique effective
divisor \(E\) such that
\[
    E\sim D.
\]
It is given by
\[
    E=D+\operatorname{div}(s),
\]
where \(s\) is viewed as a rational function after choosing the standard
rational section of \(\OO_X(D)\).
\end{proposition}

\begin{proof}
The condition \(s\in H^0(X,\OO_X(D))\) means
\[
    D+\operatorname{div}(s)\ge0.
\]
Thus \(E=D+\operatorname{div}(s)\) is effective and linearly equivalent to
\(D\).  Conversely, if \(E\ge0\) and \(E\sim D\), then \(E-D=\operatorname{div}(s)\)
for some rational function \(s\), and the inequality \(E\ge0\) says exactly
that \(s\) is a global section of \(\OO_X(D)\).  Multiplying \(s\) by a scalar
does not change \(E\), and these are the only ambiguities.
\end{proof}

\begin{proposition}
Let \(\LL\) be an invertible sheaf on a scheme \(X\), generated by global
sections \(s_0,\ldots,s_n\).  Then these sections define a morphism
\[
    \phi:X\to\PP^n
\]
such that
\[
    \LL\simeq \phi^*\OO_{\PP^n}(1).
\]
\end{proposition}

\begin{proof}
The construction of the morphism was given in
Exercise~\ref{ex:global-generation-map}.  On the open set \(U_i=\{s_i\ne0\}\), the map
lands in \(D_+(x_i)\) and is described by the regular functions \(s_j/s_i\).
The pullback of \(\OO(1)\) is trivialized on \(U_i\) by \(x_i\), and under the
map this trivialization pulls back to the trivialization of \(\LL\) by \(s_i\).
The transition functions agree, hence \(\LL\simeq\phi^*\OO(1)\).
\end{proof}

\begin{definition}
A divisor \(D\) on a smooth projective curve is \emph{base-point free} if for
every closed point \(P\), there exists a section
\[
    s\in H^0(X,\OO_X(D))
\]
with \(s(P)\ne0\) in the fiber of \(\OO_X(D)\) at \(P\).  It is \emph{very
ample} if the associated morphism to projective space is a closed immersion.
\end{definition}

\section{Degree of Line Bundles}

\begin{definition}
For an invertible sheaf \(\LL\) on a smooth projective curve \(X\), define
\[
    \deg\LL=\deg D
\]
for any divisor \(D\) with \(\LL\simeq\OO_X(D)\).
\end{definition}

\begin{proposition}
The degree of an invertible sheaf is well defined, and
\[
    \deg(\LL\tensor\MM)=\deg\LL+\deg\MM.
\]
\end{proposition}

\begin{proof}
If \(\OO_X(D)\simeq\OO_X(E)\), then \(D-E\) is principal.  Principal divisors
have degree zero, so \(\deg D=\deg E\).  If
\(\LL\simeq\OO_X(D)\) and \(\MM\simeq\OO_X(E)\), then
\[
    \LL\tensor\MM\simeq\OO_X(D+E),
\]
and the degree formula follows from additivity of degree for divisors.
\end{proof}

\section{Exercises Used Later}

\begin{exercise}[Negative degree has no sections]\label{ex:negative-degree}
Let \(X\) be a smooth projective curve and let \(D\) be a divisor with
\(\deg D<0\).  Show that \(H^0(X,\OO_X(D))=0\).
\end{exercise}

\begin{solution}
If \(0\ne s\in H^0(X,\OO_X(D))\), then
\[
    D+\operatorname{div}(s)\ge0.
\]
Taking degrees gives
\[
    \deg D+\deg\operatorname{div}(s)\ge0.
\]
Principal divisors have degree zero, so \(\deg D\ge0\), contradiction.  Hence
no non-zero section exists.
\end{solution}

\begin{exercise}[Linearly equivalent divisors]\label{ex:linearly-equivalent}
Let \(D\sim E\) on a smooth projective curve.  Show that
\[
    H^0(X,\OO_X(D))\simeq H^0(X,\OO_X(E)).
\]
\end{exercise}

\begin{solution}
Write \(E-D=\operatorname{div}(f)\).  Multiplication by \(f\) gives an
isomorphism \(\OO_X(D)\to\OO_X(E)\), because locally it shifts the allowed pole
order by exactly the divisor of \(f\).  Taking global sections gives the
desired isomorphism.
\end{solution}

\begin{exercise}[Base points and subtraction]\label{ex:basepoint-subtraction}
Let \(D\) be a divisor on a smooth projective curve and \(P\) a closed point.
Show that \(P\) is a base point of \(|D|\) if and only if
\[
    H^0(X,\OO_X(D-P))=H^0(X,\OO_X(D))
\]
inside \(H^0(X,\OO_X(D))\).
\end{exercise}

\begin{solution}
The inclusion \(\OO_X(D-P)\subset\OO_X(D)\) identifies sections of
\(\OO_X(D-P)\) with sections of \(\OO_X(D)\) that vanish at \(P\).  Thus all
sections vanish at \(P\) exactly when the two spaces of sections are equal.
This is precisely the statement that \(P\) is a base point.
\end{solution}

\begin{exercise}[Pullback preserves principal divisors]\label{ex:pullback-principal}
Let \(f:X\to Y\) be a dominant morphism of integral regular schemes and let
\(g\in k(Y)^\times\).  Show that
\[
    f^*\operatorname{div}(g)=\operatorname{div}(f^\#g)
\]
whenever the pullback divisor is defined.
\end{exercise}

\begin{solution}
At a codimension-one point \(P\in X\), the coefficient of the pullback is the
valuation of \(f^\#g\) in the DVR \(\OO_{X,P}\).  This is also the coefficient
of \(\operatorname{div}(f^\#g)\).  Since equality of Weil divisors is equality
of all codimension-one coefficients, the formula follows.
\end{solution}

\section{Divisor Class Groups on Normal Schemes}

\begin{definition}
Let \(X\) be a normal noetherian integral scheme.  The \emph{Cartier class
group} is the quotient of Cartier divisors by principal Cartier divisors.  The
\emph{Weil class group} is \(\Cl(X)\).
\end{definition}

\begin{proposition}
There is a natural injective homomorphism from Cartier divisor classes to Weil
divisor classes on a normal noetherian integral scheme:
\[
    \CaDiv(X)/\operatorname{Prin}(X)\hookrightarrow \Cl(X).
\]
\end{proposition}

\begin{proof}
A Cartier divisor is represented locally by rational functions \(f_i\).  At a
codimension-one point \(\eta\), the local ring \(\OO_{X,\eta}\) is a DVR
because \(X\) is normal and noetherian, so we may define the coefficient
\(\ord_\eta(f_i)\).  This is independent of \(i\) because transition quotients
are units.  Principal Cartier divisors go to principal Weil divisors.

If a Cartier divisor maps to a principal Weil divisor \(\operatorname{div}(g)\),
then locally \(f_i/g\) has valuation zero at every codimension-one point.  A
normal noetherian domain is the intersection of its height-one localizations
inside its fraction field.  Hence \(f_i/g\) is a unit locally, so the Cartier
divisor differs from \(\operatorname{div}(g)\) by zero.  This proves
injectivity on classes.
\end{proof}

\begin{proposition}
If \(X\) is regular, every Weil divisor is Cartier.  Hence
\[
    \Pic(X)\simeq \Cl(X)
\]
when \(X\) is regular, noetherian, integral, and every invertible sheaf is
represented by a Cartier divisor.
\end{proposition}

\begin{proof}
Let \(D=\sum n_\eta[\eta]\) be a Weil divisor.  At a codimension-one point
\(\eta\), the regular local ring \(\OO_{X,\eta}\) is a DVR, and we can choose a
local parameter \(t_\eta\).  Near \(\eta\), the local equation
\(t_\eta^{n_\eta}\) gives the prescribed coefficient.  Away from the support of
\(D\), use the equation \(1\).  Since the support is finite on affine opens,
these local equations define a Cartier divisor.  The Picard and class group
identification follows from the Cartier-Weil identification and the map
\(D\mapsto\OO_X(D)\).
\end{proof}

\section{Moving and Separating Divisors on Curves}

\begin{proposition}
Let \(X\) be a smooth projective curve and \(D\) a divisor.  The linear system
\(|D|\) separates two closed points \(P,Q\) if and only if
\[
    \ell(D-P-Q)=\ell(D)-\deg P-\deg Q.
\]
For \(k\)-rational points this says \(\ell(D-P-Q)=\ell(D)-2\).
\end{proposition}

\begin{proof}
Consider the evaluation map
\[
    H^0(X,\OO_X(D))\longrightarrow
    \OO_X(D)\tensor(\kappa(P)\oplus\kappa(Q)).
\]
Its kernel is \(H^0(X,\OO_X(D-P-Q))\).  The system separates \(P\) and \(Q\)
exactly when this evaluation map is surjective.  The target has dimension
\(\deg P+\deg Q\) over \(k\).  Therefore surjectivity is equivalent to the
displayed dimension formula.
\end{proof}

\begin{proposition}
The linear system \(|D|\) separates tangent directions at a \(k\)-rational point
\(P\) if and only if
\[
    \ell(D-2P)=\ell(D)-2.
\]
\end{proposition}

\begin{proof}
The first infinitesimal neighbourhood of \(P\) has structure sheaf
\(\OO_{X,P}/\mathfrak m_P^2\), a two-dimensional \(k\)-vector space.  The
evaluation map to this first neighbourhood has kernel
\[
    H^0(X,\OO_X(D-2P)).
\]
Separating tangent directions is exactly surjectivity of this map, and the
dimension of the target is \(2\).  The formula follows.
\end{proof}
